\section{Homework Week 13}

\begin{exercise}
    Let \(G\) be a finite group and \(p\) a prime. Prove that each element \(x\in G\)
    can be uniquely written as
    \[
        x = st,
    \]
    where \(s\) is \(p\)-regular, \(t\) is \(p\)-singular, and \(st=ts\). If
    \(x_1=s_1t_1\) is such a decomposition of an element \(x_1\) that is conjugate to
    \(x\), prove that \(s\) is conjugate to \(s_1\) and \(t\) is conjugate to \(t_1\).
\end{exercise}

\begin{proof}
    Fix \(x\in G\). Since \(G\) is finite, \(x\) has finite order. Write
    \[
        |x|=p^a m,
    \]
    where \(a\geq 0\) and \(p\nmid m\). Because \(\gcd(p^a,m)=1\), there exist integers
    \(u,v\) such that
    \[
        up^a+vm=1.
    \]
    Define
    \[
        t:=x^{vm},
        \qquad
        s:=x^{up^a}.
    \]
    Then \(s,t\in \langle x\rangle\), so \(st=ts\), and
    \[
        st
        =
        x^{up^a+vm}
        =
        x.
    \]

    We check that \(t\) is \(p\)-singular and \(s\) is \(p\)-regular. Since
    \[
        t^{p^a}
        =
        x^{vmp^a}
        =
        x^{v|x|}
        =
        1,
    \]
    the order of \(t\) divides \(p^a\), hence \(\operatorname{ord}(t)=p^k\) for some
    \(k\leq a\). Thus \(t\) is \(p\)-singular. Likewise,
    \[
        s^m
        =
        x^{up^a m}
        =
        x^{u|x|}
        =
        1,
    \]
    so \(\operatorname{ord}(s)\mid m\). Since \(p\nmid m\), we have \(p\nmid
    \operatorname{ord}(s)\), and therefore \(s\) is \(p\)-regular. This proves existence.

    For uniqueness, suppose
    \[
        x=s_1t_1=s_2t_2,
    \]
    where each \(s_i\) is \(p\)-regular, each \(t_i\) is \(p\)-singular, and
    \(s_it_i=t_is_i\). Since \(p\nmid \operatorname{ord}(s_i)\) and
    \(\operatorname{ord}(t_i)\) is a power of \(p\), we have
    \[
        \gcd\bigl(\operatorname{ord}(s_i),\operatorname{ord}(t_i)\bigr)=1.
    \]
    Hence
    \[
        |x|
        =
        \operatorname{ord}(s_i)\operatorname{ord}(t_i).
    \]
    The subgroup \(\langle s_i,t_i\rangle\) is abelian and has order
    \(\operatorname{ord}(s_i)\operatorname{ord}(t_i)=|x|\), so it equals
    \(\langle x\rangle\). In particular, \(s_i,t_i\in \langle x\rangle\).

    The cyclic group \(\langle x\rangle\) has order \(|x|=p^a m\). In a cyclic group of
    this order, an element is \(p\)-regular if and only if its order divides \(m\), and
    \(p\)-singular if and only if its order divides \(p^a\). Under the identification
    \[
        \langle x\rangle \cong \mathbb{Z}/|x|\mathbb{Z},
    \]
    the Chinese remainder theorem gives a unique decomposition of each element into
    commuting \(p\)-regular and \(p\)-singular parts. Therefore \(s_1=s_2\) and
    \(t_1=t_2\).

    Now suppose that \(x_1\) is conjugate to \(x\), say
    \[
        x_1=gxg^{-1},
    \]
    and that \(x=st\) and \(x_1=s_1t_1\) are the corresponding decompositions. Define
    \[
        s':=gsg^{-1},
        \qquad
        t':=gtg^{-1}.
    \]
    Conjugation preserves order, so \(s'\) is \(p\)-regular and \(t'\) is
    \(p\)-singular. Moreover,
    \[
        s't'
        =
        g(st)g^{-1}
        =
        g(ts)g^{-1}
        =
        t's',
    \]
    and
    \[
        s't'
        =
        g(st)g^{-1}
        =
        gxg^{-1}
        =
        x_1.
    \]
    By uniqueness of the decomposition of \(x_1\), we have \(s_1=s'\) and
    \(t_1=t'\). Hence \(s\) is conjugate to \(s_1\), and \(t\) is conjugate to \(t_1\).
\end{proof}
